%! AP Statistics — Unit 1, Lesson 1.1 — Exit Ticket KEY
\documentclass[10pt]{article}
\usepackage[margin=0.7in, top=0.55in, bottom=0.7in]{geometry}
\usepackage{xcolor}
\usepackage[most]{tcolorbox}
\usepackage{enumitem}
\usepackage{array}
\usepackage{tabularx}
\usepackage{tikz}
\tcbuselibrary{skins}

% ── Colors ───────────────────────────────────────────────────────────────────
\definecolor{navy}{HTML}{1F3A5F}
\definecolor{sky}{HTML}{EAF3FB}
\definecolor{skymid}{HTML}{C5DFF5}
\definecolor{redbg}{HTML}{FCECEC}
\definecolor{redacc}{HTML}{C0392B}
\definecolor{greenbg}{HTML}{EDF8F0}
\definecolor{greenacc}{HTML}{2D7A4F}
\definecolor{slate}{HTML}{64748B}
\definecolor{linegray}{HTML}{AAAAAA}
\definecolor{keyred}{HTML}{C0392B}

% ── Fill-in helpers ──────────────────────────────────────────────────────────
\newcommand{\ans}[1]{\textcolor{keyred}{\textit{#1}}}
\newcommand{\blank}[1]{\underline{\hspace{#1}}}

\newtcolorbox{teachernote}{
    enhanced,
    colback=redbg, colframe=redacc, boxrule=0.7pt, arc=1.5mm,
    left=2mm, right=2mm, top=1.5mm, bottom=1.5mm,
    fonttitle=\bfseries\color{white}, title={Teacher Note},
    attach boxed title to top left={yshift=-2mm, xshift=4mm},
    boxed title style={colback=redacc, colframe=redacc, arc=1mm}
}

\setlength{\parindent}{0pt}
\setlength{\parskip}{4pt}

%=============================================================================
\begin{document}

% ── Header ───────────────────────────────────────────────────────────────────
    \noindent
    \colorbox{navy}{%
        \parbox{\dimexpr\linewidth}{%
            \vspace{6pt}
            \color{white}\textbf{\large AP Statistics \hfill Unit 1, Lesson 1.1}\\[2pt]
            {\color{skymid}\small\textbf{Exit Ticket}}
            \vspace{6pt}
        }%
    }
    \vspace{1.5em}

\noindent
\begin{tabularx}{\linewidth}{X X X}
Name:\;\textcolor{keyred}{\textit{Key}} & Date:\; & Period:\;
\end{tabularx}

\vspace{0.22in}

% ── SCENARIO ─────────────────────────────────────────────────────────────────
\begin{tcolorbox}[colback=sky, colframe=navy, boxrule=0.9pt, arc=2mm,
    left=4mm, right=4mm, top=3mm, bottom=3mm]
\small\textit{A hospital tracks the recovery times (in days) of 50 patients
following knee-replacement surgery, in order to learn about typical recovery
for all patients who undergo this procedure.}
\end{tcolorbox}

\vspace{0.18in}

\textbf{For the scenario above, identify each of the following.}

\vspace{0.14in}

\renewcommand{\arraystretch}{2.2}
\begin{tabularx}{\linewidth}{@{} >{\bfseries\color{navy}}p{0.24\linewidth} X @{}}
Individual:  & \ans{A patient who has undergone knee-replacement surgery} \\
Population:  & \ans{All patients who undergo knee-replacement surgery} \\
Sample:      & \ans{The 50 patients tracked by the hospital} \\
Variable:    & \ans{Recovery time, measured in days} \\
\end{tabularx}

\vspace{0.18in}

\textbf{Does this variable vary across individuals? How do you know?}\\[6pt]
\ans{Yes — different patients recover at different rates depending on age, health,
physical therapy, and other factors. If recovery time were identical for everyone,
there would be nothing to study.}

\vspace{0.2in}

\begin{teachernote}
\scriptsize
\textbf{Scoring (completeness, not correctness):} Award credit for any reasonable, good-faith
attempt at each component. Common gaps to note for re-teaching:\\[3pt]
\textbf{Individual vs.\ Population confusion:} Students who write ``all patients'' as the
individual haven't internalized that the individual is the \emph{unit of observation},
not the group. Ask: ``What does each row in the data table represent?''\\[3pt]
\textbf{Sample vs.\ Population:} Some students write the same thing for both. Push them:
``Did the hospital study \emph{every} patient who has ever had this surgery?''\\[3pt]
\textbf{Variable:} Accept ``recovery time'' without units. Flag answers that list the
\emph{individual} again (e.g., ``the patients'') — this indicates the student doesn't yet
distinguish between a person and a measurement.
\end{teachernote}

\vspace{0.18in}

% ── AP-STYLE MC ──────────────────────────────────────────────────────────────
\begin{tcolorbox}[colback=redbg, colframe=redacc, boxrule=0.8pt, arc=2mm,
    left=4mm, right=4mm, top=3mm, bottom=3mm]
\small\textbf{AP-Style Practice}\\[6pt]
A researcher surveys 200 randomly chosen adults about their daily exercise habits.
Which of the following best describes the \textbf{population}?

\vspace{10pt}
\begin{enumerate}[label=\textbf{(\Alph*)}, leftmargin=2em, itemsep=6pt]
  \item The 200 adults who were surveyed
  \item The number of minutes of exercise reported each day
  \item \textbf{\textcolor{keyred}{All adults}} \textcolor{keyred}{\textit{$\leftarrow$ correct}}
  \item The researcher conducting the survey
\end{enumerate}
\end{tcolorbox}

\vspace{0.1in}
\noindent\textcolor{keyred}{\textbf{Answer: C}} \quad
\ans{The population is the full group the researcher wants to draw conclusions
about — all adults. The 200 surveyed are only the sample. (B) is the variable;
(D) is not part of the data at all.}

\vspace{0.12in}
\begin{teachernote}
\textbf{Most common wrong answer: A.} Students who choose A are confusing the
sample (who was studied) with the population (who we want to know about).
This is the single most important conceptual distinction to nail in Lesson 1.1.
If more than a third of the class chooses A, revisit the population vs.\ sample
distinction explicitly at the start of Lesson 1.2.
\end{teachernote}

\end{document}
